% ============================================================================
%  spiral_v15_preprint.tex  (v16-arxiv-ready)
%
%  A per-layer negative-control ledger for a multi-stage simulation pipeline.
%
%  Framing (must hold): this is a NUMERICAL-EXPERIMENT / METHODOLOGY report,
%  not a theory of everything, and not a proof that the seven-layer framework
%  is established. Every number is transcribed from the frozen reference run
%  (1245 s), generated by running spiral_model_v15.py with spiral_metric_v15.py.
%  Source files required: spiral_model_v15.py, spiral_metric_v15.py only.
%  All outputs (_v15_run.log, result/_v15_data.json, result/spiral_v15.png,
%  result/spiral_v15_metrics.png)
%  are generated from scratch; no pre-existing data is read.
%
% Compile-tested: pdflatex → bibtex → pdflatex ×2, passes cleanly.
% ============================================================================
\documentclass[11pt,a4paper]{article}

\usepackage[numbers,sort&compress]{natbib} %
\usepackage[utf8]{inputenc}
\usepackage[T1]{fontenc}
\usepackage{amsmath,amssymb}
\usepackage{booktabs}
\usepackage{graphicx}
\usepackage{array}
\usepackage[margin=2.4cm]{geometry}
\usepackage[hidelinks]{hyperref}
\usepackage{enumitem}
\usepackage{xcolor}

\newcommand{\code}[1]{\texttt{\small #1}}
\newcommand{\good}[1]{\textcolor{black}{\textbf{#1}}}

\title{\bfseries A negative-control ledger  \\[2pt]
for multi-stage numerical pipelines \\[8pt]
\large --- demonstrated on Ising/MERA\\}

\author{%
   JianZhong Yan\\
  \small \textit{(independent researcher)}\\
  \small \texttt{mohedanovinita984@gmail.com}\\
}
\date{\today}

\begin{document}
\maketitle

% ---------------------------------------------------------------------------
\begin{abstract}
\noindent
We state an operational criterion for a genuine \emph{negative control} in
multi-stage numerical pipelines: \textbf{it must be able to fail when the
conclusion does not hold}. Applying this criterion excludes positivity,
finiteness, and non-emptiness checks, and sorts surviving controls into four types
(repulsion, null-model, two-path, upper-bound). The ledger built from them carries
a \emph{completeness guard} that checks cited guards against the runtime registry,
making the ledger hard to fake. Checks whose guarded proposition is measurably
false are registered as \emph{diagnostics} and excluded from the pass-rate
denominator.

We demonstrate this method on a seven-layer Ising/MERA pipeline.
\textbf{Results}: 54 of 54 scored assertions hold, 18 of 18 metrics meet target,
and 12 negative controls cover 6 of 7 layers. The pass rate is a statement about
guard execution, \emph{not} about the truth of any physical emergence claim. The
substrate is pinned by an external anchor---the ground-state energy density lies
on the parameter-free CFT finite-size formula $e_\infty-\pi cv/(6L^2)$ \emph{at
every} $L$ (relative residual $1.81\times10^{-6}$ at $L{=}16$). A low-energy
spectral ratio supplies a second ruler for the central charge
($0.212\%$ deviation), which is \emph{independent of the entanglement-spectrum
family} (though we note explicitly that the $0.212\%$ figure and the
``ratio $\to 8$'' statement are the \emph{same measurement}, not a double
independent check). We report three \emph{negative} results: one layer has no
causal path to any downstream consumer; the bulk-geometry readout in the current
fixed-wiring MERA implementation does not depend on the fitted state (it is an
analytic property of the binary network layout); and two self-contradictory
reporting conventions were corrected.

\smallskip
\noindent\textbf{Boundary.} ``$100\%$'' means 54 scored assertions hold. It is
\emph{not} evidence that the system under study became better understood, and this
paper claims no physical result beyond textbook Ising criticality.
\textbf{Code availability:} \url{https://github.com/yanjianzhong/spiral-emergence}
\end{abstract}

% ---------------------------------------------------------------------------
\section{Introduction}
\label{sec:intro}

\subsection{The problem: a sanity check is not a negative control}

A layered pipeline derives quantities stage by stage. At each stage an assertion is
made --- ``the spectrum is degenerate'', ``the entropy is positive'', ``the grid is
fine enough'' --- and downstream stages consume the derived quantity. In practice
these assertions are enforced by \emph{positivity}, \emph{finiteness}, and
\emph{non-emptiness} checks, often promoted from a hard \code{raise} into a
readable boolean.

The difficulty is structural rather than clerical: such checks are \emph{consistent
with any conclusion}. ``Schmidt rank $\ge 2$'' holds for a correct calculation and
for a calculation that is wrong in every way that matters. A pipeline can therefore
report a wall of green checks while every substantive claim in it is false, and
nothing in the report distinguishes the two situations.

This paper is about a discipline that fixes that, and about what happens when the
discipline is applied honestly --- including where it \emph{cannot} be applied.

\subsection{Contributions}

\begin{enumerate}[leftmargin=1.4em,itemsep=2pt]
  \item \textbf{A criterion} (Sec.~\ref{sec:criterion}) that separates genuine
        negative controls from true-by-construction checks, and that explicitly
        \emph{excludes} the positivity/finiteness/non-emptiness family.
  \item \textbf{A taxonomy} (Sec.~\ref{sec:taxonomy}) of four control types, each
        tied to a concrete implemented instance.
  \item \textbf{A ledger with a completeness guard} (Sec.~\ref{sec:ledger}) that
        checks the ledger against the \emph{runtime} registry of guards, rather
        than trusting the ledger's own self-report. Crucially, this guard is
        itself falsifiable: if a cited guard does not exist or a layer is
        missing, it fails --- making the ledger \emph{hard to fake}
        (Sec.~\ref{sec:repro}).
  \item \textbf{A scoring/diagnostic split} (Sec.~\ref{sec:split}): checks whose
        guarded proposition is measurably false do not enter the denominator.
  \item \textbf{A demonstration} (Secs.~\ref{sec:pipeline}--\ref{sec:results}) on a
        seven-layer pipeline, reported together with its gaps, its negative
        results, and one real defect that the controls caught
        (Sec.~\ref{sec:bug}).
\end{enumerate}

\subsection{Scope, and what this paper does not claim}
\label{sec:nonclaims}

Stated early, because it governs how every number below should be read:

\begin{itemize}[leftmargin=1.4em,itemsep=2pt]
  \item We do \textbf{not} claim that the seven-layer framework under study is
        established. Its independently audited maturity is A\,$\times$\,1,
        B\,$\times$\,3, C\,$\times$\,3, D\,$\times$\,0 (Sec.~\ref{sec:gaps}).
  \item We do \textbf{not} claim that any holographic bridge was established
        between layers 4 and 5 (Sec.~\ref{sec:l4l5}); the tensor-network/holographic
        correspondence is considered only as background motivation following
        Swingle's entanglement-renormalization approach to emergent spacetime
        \cite{Swingle2012}.
  \item We do \textbf{not} claim that a $100\%$ pass rate means the evidence
        became stronger (Sec.~\ref{sec:diagnostics}).
  \item The physics used as a substrate --- critical Ising/CFT scaling --- is
        \textbf{textbook material}, verified here as a substrate, not offered as a
        result.
  \item We do \textbf{not} claim that MERA cannot produce state-dependent geometry
        in general; we show only that in the present fixed-wiring binary
        implementation, the readout is an analytic property of the graph layout
        (Sec.~\ref{sec:l5geometry}).
\end{itemize}

The contribution on offer is methodological and engineering: the substrate is
pinned down by external anchors, and the novelty, such as it is, lives in the
\emph{method layer}.

% ---------------------------------------------------------------------------
\section{Method: criterion, taxonomy, and ledger}
\label{sec:method}

\subsection{The criterion}
\label{sec:criterion}

\begin{quote}
\textbf{Criterion.} A check qualifies as a negative control only if there exists a
concrete, reachable state of the pipeline in which \emph{the conclusion it guards
is false} and the check \emph{fails}.
\end{quote}

The criterion is applied as a filter, not as an aspiration. In the pipeline studied
here, it removed an existing candidate --- ``the entanglement entropy is positive''
at layer 4 --- which had been counted as a control. Positivity holds for a wrong
entropy as readily as for a right one; it is a sanity check and is now labelled as
one. The same filter excludes the whole positivity / finiteness / non-emptiness
family. This exclusion is recorded in the table header of the control registry in
the source (\code{spiral\_metric\_v15.\_NEG\_CTRL\_TABLE}), so that the exclusion
is a property of the code rather than of this prose.

\subsection{Four types of control}
\label{sec:taxonomy}

Every control admitted by the criterion is assigned one of four types. Each type is
defined together with the instance that actually implements it in the pipeline:

\begin{table}[h]
\centering\small
\begin{tabular}{@{}p{2.1cm}p{5.6cm}p{7.2cm}@{}}
\toprule
\textbf{Type} & \textbf{Definition} & \textbf{Implemented instance} \\
\midrule
Repulsion &
A quantity is pushed away from a value it would take if the conclusion were false;
the observed displacement is the evidence. &
Layer 4 bond dimension: when the bond dimension is forced below its lower bound,
the downstream quantity degrades by $6.47\times$ (criterion $>1.5\times$). \\[3pt]
Null model &
The measurement is compared against an ensemble generated to share the
\emph{nuisance} structure of the real object while destroying the structure under
test. &
Layer 5 bulk geometry: three null families $\times$ 3 instances --- \code{gnm}
(matched $|V|$, $|E|$; controls scale), \code{ws} (degree-matched small world;
controls degree distribution), \code{tree} (balanced binary tree; controls the
``nearly a tree'' premise). Two-sided acceptance band $[5\%,95\%]$. \\[3pt]
Two-path &
The same quantity is obtained by a second, structurally different route; agreement
is the evidence. &
Five instances, spread across four layers (L2 $\times$2, L3, L5, L6) --- for example
the L2 spectral-gap degeneracy guard. \\[3pt]
Upper bound &
A computed quantity must not exceed an independently known bound. &
Layer 7: a single elementary bound. This is the thinnest control in the ledger, and
is reported as such. \\
\bottomrule
\end{tabular}
\caption{The four control types, with the instance that implements each. Type
counts in the demonstrated ledger: repulsion 5, two-path 5, null-model 1,
upper-bound 1 (12 total).}
\label{tab:types}
\end{table}

\subsection{The ledger and its completeness guard}
\label{sec:ledger}

The ledger groups controls by physical layer, and each entry cites the guard name(s)
that implement it. A ledger is a self-report, and a self-report can be written as a
brochure: it can cite guards that do not exist, or quietly omit a layer.

The completeness guard closes that hole by checking the ledger against the
\emph{runtime registry} of guards rather than against its own text:

\begin{itemize}[leftmargin=1.4em,itemsep=2pt]
  \item every guard name cited by the ledger must exist in the live registry;
  \item each of the seven layers must appear exactly once.
\end{itemize}

\noindent In the reference run this guard reports \code{phantom references: none;
missing layers: none}. Section~\ref{sec:bug} describes the defect this guard caught.

\begin{quote}
\textbf{Note on what this guard is not.} It is not a physical conclusion. It
guarantees only that the ledger contains no phantom citation and no silently
dropped layer. Its object is the honesty of the accounting, not the truth of the
physics.
\end{quote}

\medskip
\noindent\textbf{Why this guard is the method's sharpest edge.}
The completeness guard is not merely a bookkeeping aid; it is the mechanism that
makes the \emph{method itself} falsifiable. A ledger that cites phantom guards or
silently omits layers is, in effect, a falsified document --- and this guard is
designed to detect exactly that state. This is what distinguishes the present
approach from every other pipeline-validation framework we are aware of: the
validation framework is subject to the same logical standard it imposes on the
pipeline. If the ledger's own accounting is wrong, the guard fails
\emph{first}, before any physical claim is evaluated. The defect caught in
Sec.~\ref{sec:bug} --- where the completeness criterion itself was inverted ---
is the proof that this self-falsifiability is not a slogan: it caught a bug
\emph{in the validation layer}, which is exactly where bugs are hardest to see.

\subsection{Scoring items versus diagnostics}
\label{sec:split}

Some guards protect propositions that are \emph{measurably false} in the system as
built, or that depend on a parameter rather than on the system. Reporting these as
failures would be misleading, and silently dropping them would be worse. They are
therefore registered as \emph{diagnostics} and excluded from the pass-rate
denominator.

In the reference run: 74 guards total, of which 54 are scored and 20 are
diagnostics. \textbf{Of the 20 diagnostics, 17 fail by design.} The
\code{expect\_pass} flag is set explicitly at registration, so the distinction is a
property of the code and survives any re-reading of the output. The rationale for
each of the four most recent additions is recorded individually
(Sec.~\ref{sec:diagnostics}).

\medskip
\noindent\textbf{Why this matters.} A pipeline that counts ``checks that hold'' as
evidence will inflate its score by admitting checks that cannot fail. Separating
the denominator is the cheapest available defence, and it costs one boolean per
guard.

% ---------------------------------------------------------------------------
\section{The demonstration pipeline}
\label{sec:pipeline}

\subsection{Seven layers}

The pipeline derives, layer by layer, from an exact-diagonalisation (ED) solution
of the critical transverse-field Ising chain through entanglement spectra, bond
dimensions, tensor-network bulk geometry, and a downstream scaling analysis. Each
layer exposes (i) the proposition it asserts, (ii) the derived quantity it hands
downstream, and (iii) its downstream consumers. The audit table records, per layer,
which downstream consumer actually reads the derived quantity --- this is what makes
the gap in Sec.~\ref{sec:l1gap} visible.

\subsection{Substrate and its constraint}
\label{sec:substrate}

The validation substrate is the critical transverse-field Ising chain: $c = 1/2$,
$v = 2$, $e_\infty = -4/\pi$. Landau--de~Gennes (Gray--Scott) pattern formation
appears as a second, non-critical exhibit used for the layer-3 descriptor analysis.

One boundary deserves emphasis because it is easy to misattribute. The pipeline's
chain-length ceiling ($L \le 16$ on the closed loop) comes from a
\textbf{tensor-network library constraint}: the MERA construction used \cite{EvenblyVidal2009} requires $L$
to be a power of two. It is \emph{not} an exact-diagonalisation limit (sparse ED
reaches $L = 20$ in $22.4$~s) and \emph{not} a laptop budget. This cause was
mis-stated twice in the source comments before being corrected; it is reported here
because the same misattribution is easy to make elsewhere.

% ---------------------------------------------------------------------------
\section{Results}
\label{sec:results}

\subsection{Coverage of the ledger, and its gap}
\label{sec:gaps}

The ledger's 12 negative controls cover \textbf{6 of the 7 layers}; the gap is
\code{['L1']}. The strongest layer is L5 (three null families $\times$ 3 instances,
two-sided band); the thinnest is L7, with a single upper-bound control. The
per-layer coverage guard is \emph{deliberately not scored}: coverage is a to-do
list, not a property of the model, and scoring it would be grading the checklist.

\begin{table}[h]
\centering\small
\begin{tabular}{@{}lrl@{}}
\toprule
\textbf{Quantity} & \textbf{Reference run} & \textbf{Note} \\
\midrule
Scored assertions & $54/54 = 100.00\%$ & see Sec.~\ref{sec:diagnostics} \\
Diagnostics excluded & 20 (17 failing by design) & \code{expect\_pass=False} \\
Metrics meeting target & $18/18$ & \\
Negative controls & 12 & repulsion 5 / two-path 5 / null 1 / bound 1 \\
Layer coverage & $6/7$, gap \code{['L1']} & Sec.~\ref{sec:l1gap} \\
Ledger completeness & no phantom refs, no missing layer & Sec.~\ref{sec:ledger} \\
Wall-clock (single laptop) & $1245$~s & order-of-magnitude only; was $1138$~s in an earlier run \\
\bottomrule
\end{tabular}
\caption{Reference-run summary. Values transcribed from the pipeline log of the
reference run.}
\label{tab:summary}
\end{table}

\subsection{External anchor: the $1/L^2$ coefficient, at every $L$}
\label{sec:anchor}

The substrate is pinned by a comparison to a formula with \emph{no free parameters} \cite{CalabreseCardy2004}:
\begin{equation}
  \frac{E_0(L)}{L} \;=\; e_\infty \;-\; \frac{\pi c v}{6 L^2},
  \qquad c = \tfrac12,\quad v = 2,\quad e_\infty = -\tfrac{4}{\pi}.
  \label{eq:cft}
\end{equation}
$E_0$ comes from exact diagonalisation; the constants on the right do not enter its
computation.

\begin{table}[h]
\centering\small
\begin{tabular}{@{}rll@{}}
\toprule
$L$ & relative residual vs.\ Eq.~\eqref{eq:cft} & residual-inverted power $p$ \\
\midrule
8  & $2.88\times10^{-5}$ & 4.007 \\
$\vdots$ & $\vdots$ & $\vdots$ \\
16 & $\mathbf{1.81\times10^{-6}}$ & 4.004 \\
18 & $1.13\times10^{-6}$ & 4.002 \\
\bottomrule
\end{tabular}
\caption{The residual decreases monotonically with $L$, and the residual-inverted
power is $p = 4.00$ throughout --- the expected leading correction. The anchor
script carries a \code{-{}-selftest} that demonstrates the comparison is not
identically true.}
\label{tab:anchor}
\end{table}

\noindent\textbf{Honest boundary.} This is stronger than the anchor already present
in the code, which compared only the \emph{extrapolated} $e_\infty$ against
$-4/\pi$ (deviation $5.0\times10^{-6}$); here the comparison is against the
\emph{per-$L$} $1/L^2$ coefficient, which involves the product $cv$. But it remains
a statement that \emph{exact diagonalisation was performed correctly}. It is
textbook content, not a result of this paper.

\subsection{A second family that is genuinely independent}
\label{sec:secondfamily}

A second estimator for the central charge uses the low-energy spectral ratio
$(E_2 - E_0)/(E_1 - E_0)$, which is independent of $v$:

\begin{table}[h]
\centering\small
\begin{tabular}{@{}lrl@{}}
\toprule
\textbf{Quantity} & \textbf{Value} & \textbf{Note} \\
\midrule
ratio at $L{=}8$  & $7.9231$ & monotone toward 8 \\
ratio at $L{=}16$ & $7.9807$ & \\
$c$ interval & $[0.49894,\,0.50034]$ & max deviation $0.212\%$ \\
Rényi family, same pipeline & --- & deviation $7.97\%$ \\
\bottomrule
\end{tabular}
\caption{Low-energy spectral estimator versus the entanglement-spectrum family.}
\label{tab:second}
\end{table}

\noindent\textbf{Two hard boundaries, both required reading.}
\emph{(i)} The estimator is a self-consistency check, not an assumption-free
measurement: the operator content is identified before use ($E_1$ is the $R$-sector
vacuum, $\Delta_1 = 1/8$, at every point; $E_2$ is the NS two-particle state,
$\Delta_2 = 1$), so $\Delta_1 = 1/8$ enters the construction.
\emph{(ii)} The readout does \textbf{not} converge monotonically in $L$ ---
$L{=}16$ is the worst point --- because the residual is an $O(1/L)$ finite-size
correction. One may say ``consistent to $0.21\%$''; one may \textbf{not} say
``extrapolates to $0.5$''.

\medskip
\noindent\textbf{A trap worth naming explicitly.} The $0.212\%$ figure and the
``ratio $\to 8$'' statement are \emph{the same measurement}: $0.212\%$ \emph{is}
the $0.24\%$ by which the ratio misses 8. It must therefore \textbf{not} be written
as ``the spectral family and the Rényi family corroborate each other''. The correct
statement is that the spectral family is \emph{independent of the
entanglement-spectrum family}. Those are different claims, and only the second is
supported.

\subsection{Qualitative distinction: area law versus logarithmic law}
\label{sec:arealaw}

At the critical point the entanglement scaling is logarithmic; away from it, the
area law holds. The measured peak of $|c_{\text{scaling}}|$ sits at $h/J = 1.00$;
at the critical point $c_{\text{scaling}} = 0.4976$ ($R^2 = 1.0000$), while the
off-critical points give $0.0007$ ($h/J = 0.5$) and $-0.0187$ ($h/J = 1.5$). A
shape test on $S(\ell)$ gives $\Delta\mathrm{AIC} = 76.14 / 10.61 / 16.35$.

\noindent\textbf{Honest boundary.} The shape fit ($c = 0.5072$) and the $L$-scaling
fit ($c = 0.4976$) use \emph{the same} $S$ data; this is self-consistency, not a
cross-check. The lever arm $L: 8 \to 16$ spans only $1.00$ octave in $\log_2 L$.
This is a \emph{qualitative} distinction, not a quantitative extrapolation.

\subsection{Negative controls caught a real defect}
\label{sec:bug}

The strongest argument that the controls are not decorative is that one of them
caught a defect in which \textbf{a completeness criterion was inverted}.

At three places the ledger wrote ``the XX negative control has been run'' --- a
\emph{self-reported} name --- into the same tuple as the actually-executed guard
names. Such names are registered only on an early-return branch, and therefore
\emph{never enter the registry when the control executes normally}. The consequence
was exactly backwards: the completeness check \textbf{failed in the healthy case and
passed when the control was missing}. Phantom references: 3 before the fix, none
after.

\medskip
\noindent\textbf{What makes this worth reporting.} The defect was found only by a
\emph{full pipeline run}. Before that, a source-text-versus-regex audit had reported
$26/26$ correct. A source-text audit is necessary but not sufficient: it cannot see
runtime registration paths. The defect was introduced by the present work, not
inherited; it is reported because the failure mode --- a check that passes for the
wrong reason --- is the exact failure this entire method exists to prevent.

% ---------------------------------------------------------------------------
\section{Honest boundaries and negative results}
\label{sec:boundaries}

This section is a first-class section, not a discussion appendix. The method's
central claim is that a pipeline should account for what it \emph{cannot} check; if
that accounting is demoted to a footnote, the method has not been applied.

\subsection{Twenty diagnostics, seventeen failing by design}
\label{sec:diagnostics}

\noindent Log line, quoted: \code{guards 54/54 = 100.00\% (plus 20 diagnostics
excluded, of which 17 did not pass)}.

``Did not pass'' is not a synonym for ``defective''. These 17 are \emph{archived
negative results}. The four most recent additions, each with its reason for
exclusion from the denominator:

\begin{itemize}[leftmargin=1.4em,itemsep=3pt]
  \item \textbf{A third corroborating path does not exist.}
  \item \textbf{The MERA bulk readout does not depend on the state}
        (Sec.~\ref{sec:l5geometry}).
  \item \textbf{That layer has no causal path} (Sec.~\ref{sec:l1gap}).
  \item \textbf{The spectral radius is parameter-dependent}, not a property of the
        system.
\end{itemize}

\noindent By contrast, the one addition that \emph{does} enter the denominator can
fail: if the bond-dimension derivation were replaced by a hand-picked constant, its
control would flip. The threshold there ($>1.5\times$) is an order-of-magnitude
comparison, and the line it references is not a strict lower bound --- both facts
are recorded in the guard's own note.

\subsection{Gap 1 --- a layer where no control can be constructed}
\label{sec:l1gap}

For layer L1 the ledger reports no negative control. The cause has been traced, and
it is not ``nobody got round to it''.

The quantity L1 hands downstream, \code{equal\_state}, is \textbf{never consumed
anywhere in the repository}: it appears exactly once, in its own \code{return}
statement. Sibling quantities are no better --- a condition number is reported but
nothing is checked against it; a dimension is fed into a downstream derivation whose
result is \emph{printed} but never compared against $2^{L_{\text{chain}}}$.

There is no quantity to push away from, because there is no downstream consumer to
be affected. The gap is registered as a \emph{source-level attribute-access scan},
which is falsifiable in the useful direction: should anyone connect the quantity, the
scan flips. A smoke test injects a fake consumer to demonstrate that the scan is not
identically false.

\medskip
\noindent\textbf{This gap is ``empty for a reason'', not ``left unfilled''.} It has
deliberately \emph{not} been entered into the L1 row of the ledger --- doing so would
pretend that L1 has a negative control.

\subsection{Gap 2 --- a bridge that is wired, and still carries nothing}
\label{sec:l4l5}

An earlier internal note claimed that the L4$\to$L5 bridge was \emph{empty}. Direct
reading of the call sites shows that claim was wrong: the bridge is wired end to
end --- the MERA graph is built, the curvature report is computed from it, the
forman value is derived from it, and the control family is constructed on it.

The controls are not missing either; they are the strongest in the ledger (three
null families $\times$ 3 instances, two-sided band). Measured outcomes:

\begin{itemize}[leftmargin=1.4em,itemsep=2pt]
  \item \textbf{F6b} (the MERA bulk geometry lies within the control distribution):
        \emph{passes}. Forman percentile $33\%$, Ollivier percentile $33\%$ inside
        the $[5\%,95\%]$ band.
  \item \textbf{F6c} (the MERA geometry is more negatively curved than \emph{all}
        controls): \emph{false, by design excluded}. Per family, Forman:
        \code{gnm} $-3.553$, \code{tree} $-0.867$, \code{ws} $-0.194$, versus MERA
        $-2.242$. The \code{gnm} control is \emph{more} negative than the MERA
        graph --- the escalation claim does not even have the right sign.
  \item \textbf{F6d} (the readout depends on the state, or on the bond dimension):
        \emph{false, by design excluded}. See Sec.~\ref{sec:l5geometry}.
\end{itemize}

\noindent So the gap is neither ``not wired'' nor ``no control available''. It is a
third kind: \textbf{no upstream-derived quantity flows along the wire.} The quantity
L4 derives ($\chi_{\text{dev}} = 4$, from the area law) is consumed elsewhere --- its
consumers were enumerated by search --- and \emph{none of them leads to the L5
construction}, which hard-codes its own bond dimension.

\subsection{The ``geometry'' readout is an analytic property of the wiring}
\label{sec:l5geometry}

Layer 5's bulk-geometry readout is computed on a MERA tensor-network graph. Within
the present implementation---a fixed-wiring binary MERA where the graph builder
reads only the index map---the readout does not depend on the fitted state, seed,
or bond dimension (Table~\ref{tab:invariance}). Three measurements show this:

\begin{table}[h]
\centering\small
\begin{tabular}{@{}lll@{}}
\toprule
\textbf{Varied} & \textbf{Method} & \textbf{Result} \\
\midrule
seed ($0,1,7$)      & rebuild and compare edge sets & \textbf{identical}, bit for bit \\
bond dimension $\chi$ ($2,4,8$) & rebuild and compare edge sets & \textbf{identical}, bit for bit \\
state               & the graph builder reads only the index map;
                      the fit routines only replace tensor data & wiring never changes \\
\bottomrule
\end{tabular}
\caption{The MERA bulk graph does not depend on the state, on the seed, or on the
bond dimension in the current fixed-wiring implementation.}
\label{tab:invariance}
\end{table}

\noindent\textbf{Honest boundary.} In homogeneous MERA tensor networks, many
structural properties are predetermined by the network geometry and are largely
independent of variational parameters \cite{EvenblyVidal2011}. This is a normal
property of the architecture; what it means here is that the readout is a property
of the \textbf{binary-layout wiring}, and the upstream claim must be re-worded
accordingly. It is therefore \emph{not} a state-dependent emergent geometry under
the present implementation. Whether a modified architecture (e.g., cMERA,
state-dependent disentanglers, or gauged/SymTFT lift) restores state dependence is
left to future work following Swingle's entanglement-renormalization approach
\cite{Swingle2012}. Reported values: $|V| = 62$, $|E| = 91$, $\bar d = 2.935$,
\textbf{triangles $= 0$}, Forman mean $-2.242$, Ollivier mean $-0.3297$,
percentiles $33\%/33\%$.

\medskip
\noindent\textbf{Distinguishing the two gaps.} They must not be merged into one
``gap'' in a report:

\begin{table}[h]
\centering\small
\begin{tabular}{@{}lll@{}}
\toprule
 & \textbf{L1} & \textbf{L4$\to$L5} \\
\midrule
path      & no causal path (unconsumed quantity) & path complete end to end \\
control   & cannot be constructed (nothing to push) & constructed; strongest in the ledger \\
what is missing & the \emph{tool} & the \emph{alignment} \\
\bottomrule
\end{tabular}
\caption{Two structurally different gaps.}
\label{tab:gaps}
\end{table}

\noindent In one sentence: \textbf{L1 is empty for a reason; L4$\to$L5 is filled
generously, but on a different question than the one the framework asks.}

\subsection{Two rulers for a length range}
\label{sec:tworulers}

When reporting a range of chain lengths, the ratio and the logarithmic span must be
reported separately. $L: 8 \to 16$ is a factor of $2.00$ but spans only
$\mathbf{1.00}$ \textbf{octave} in $\log_2 L$; $L: 8 \to 18$ is $2.25\times$ but
$1.17$ octaves. Reporting the first number while implying the second is an
overstatement. The ceiling itself comes from the tensor-network library constraint
discussed in Sec.~\ref{sec:substrate}.

\subsection{Two self-contradictory conventions, corrected}

Two reporting conventions internal to the project contradicted each other and were
corrected rather than reconciled by wording: (i) a claim that a scatter statistic
was ``inter-seed spread on the reference side'' --- contradicted by the archived
fields; (ii) a claim that two routes ``corroborate each other'' --- contradicted by
the fact that they are the same measurement (Sec.~\ref{sec:secondfamily}). Both are
recorded as corrections, with the original documentation left unedited and the
discrepancy noted.

% ---------------------------------------------------------------------------
\section{Related work}
\label{sec:related}

\textbf{Tensor networks and MERA.} The Multi-scale Entanglement Renormalization
Ansatz (MERA) was introduced by Vidal \cite{Vidal2007} and developed into a
binary-graph implementation by Evenbly and Vidal \cite{EvenblyVidal2009}. Its
application to holographic duality and emergent geometry was motivated by Swingle
\cite{Swingle2012}, who noted that entanglement renormalization naturally produces
a hyperbolic geometry in the bulk. In homogeneous MERA, many structural properties
are predetermined by the network geometry and largely independent of variational
parameters \cite{EvenblyVidal2011}, which is the context for our negative result
in Sec.~\ref{sec:l5geometry}.

\textbf{Negative controls in experiment.} The criterion used here is the ordinary
logic of a controlled experiment, transplanted into software. The transplant is the
point: in a wet-lab experiment a negative control is a physical sample; in a
pipeline it must be an \emph{executable} alternative, and executable alternatives
are easy to omit without noticing.

\textbf{Property-based testing and metamorphic testing.} The repulsion and two-path
types are close relatives of metamorphic relations and property-based tests, in that
they assert a relation that must hold across a transformation rather than a fixed
expected output. The difference is one of framing: these controls are attached to
\emph{scientific claims about a physical system}, and their failure is meant to
propagate into how loudly the claim is stated.

\textbf{Reproducibility infrastructure.} Workflow engines and provenance systems
record \emph{what ran}. This work is about a complementary question: \emph{what
would have to be true for the result to be wrong}, and whether the pipeline would
notice.

\textbf{Model validation in the physical sciences.} Null-model comparison is
standard practice in network science, and the null-model control here
(Sec.~\ref{sec:taxonomy}) is borrowed directly from that tradition --- including the
discipline of reporting \emph{which} nuisance structure each null destroys.

\textbf{What traditional sanity checks miss.} The most direct way to appreciate
the value of the negative-control discipline is to ask: which defects in this
pipeline would \emph{not} have been caught by conventional positivity,
finiteness, or non-emptiness checks? Three concrete cases answer that question:

\begin{enumerate}[leftmargin=1.4em,itemsep=3pt]
  \item \textbf{An inverted completeness criterion} (Sec.~\ref{sec:bug}): a guard
        that \emph{passed} when the referenced control was missing and
        \emph{failed} in the healthy case. Every positivity and finiteness check
        in the pipeline reported green throughout. Only the completeness guard ---
        which checks that cited guard names exist in the runtime registry ---
        could detect this, because the defect was in the \emph{accounting
        structure}, not in any numerical value.
  \item \textbf{A fully wired bridge carrying no upstream quantity}
        (Sec.~\ref{sec:l4l5}): layer 4 derives $\chi_{\text{dev}}=4$ from the
        area law, but layer 5 hard-codes its own bond dimension. The call chain
        is complete end-to-end; every function returns a finite, non-empty
        result. A positivity check on the curvature report would have passed
        silently. The gap is visible only to a control that asks whether the
        \emph{upstream-derived quantity} actually flows to the consumer.
  \item \textbf{A readout invariant under the fitted state}
        (Sec.~\ref{sec:l5geometry}): the MERA bulk graph is bit-identical across
        seeds, bond dimensions, and \emph{states}. A sanity check asking
        ``is the graph non-empty?'' or ``is curvature finite?'' passes
        trivially. Only the explicit invariance test (rebuild and compare edge
        sets) reveals that the readout is an analytic property of the wiring,
        not an emergent geometry.
\end{enumerate}

\noindent These are not hypothetical vulnerabilities. They are defects that
\emph{existed in the pipeline as built}, and that the negative-control ledger
caught precisely because each control is required to be able to fail. A
positivity check cannot catch what is already positive.

% ---------------------------------------------------------------------------
\section{Reproducibility}
\label{sec:repro}

\textbf{Source files.} The pipeline requires exactly two Python modules to run:
\code{spiral\_model\_v15.py} (entry point; orchestrates all seven layers) and
\code{spiral\_metric\_v15.py} (guard registry, ledger logic, metric definitions,
negative-control table). No ancillary data files, precomputed logs, or external
datasets are needed; every output artefact is generated from scratch.

\medskip
\noindent\textbf{Entry point.} A single command, no arguments:

\begin{quote}
\code{python spiral\_model\_v15.py}
\end{quote}

\noindent\textbf{Environment and dependencies.} The pipeline requires Python 3.x
with \code{numpy}, \code{scipy}, and a tensor-network library (the MERA
construction follows Evenbly \& Vidal \cite{EvenblyVidal2009}, requiring $L$ as a
power of two; this is the source of the chain-length ceiling, not an ED limit).
Exact library versions are pinned in \code{requirements.txt} in the repository.
On a single laptop the run takes on the order of $20$ minutes.

\medskip
\noindent\textbf{Generated outputs.} The run writes all artefacts to disk:
\code{\_v15\_run.log} (full pipeline log, $\sim$1245~s),
\code{result/\_v15\_data.json} (structured metrics), 
\code{result/spiral_v15.png} , and
\code{result/spiral_v15_metrics.png} (20-panel metric report). The guard ledger is
printed to standard output. No pre-existing data is read; the pipeline is fully
self-contained.

\medskip
\noindent\textbf{Code availability.} The complete source code,
\code{requirements.txt}, and documentation are available at:
\url{https://github.com/yanjianzhong/spiral-emergence}
(tag \texttt{v15}). No ancillary material is uploaded to arXiv; reproduction
requires only cloning the repository and running the entry point.

\medskip
\noindent\textbf{Expected output.} The following lines are what a correct run prints. They
are the reader's check that the reproduction succeeded:

\begin{table}[h]
\centering\small
\begin{tabular}{@{}p{8.4cm}p{6.6cm}@{}}
\toprule
\textbf{Line} & \textbf{Expected} \\
\midrule
\code{metrics meeting target} & \code{18/18} \\
\code{guard pass rate} & \code{54/54 = 100.00\%} (plus 20 diagnostics excluded,
                         17 of them not passing) \\
\code{ledger completeness} & \code{phantom references: none; missing layers: none} \\
negative-control layer coverage & $6/7$, gap \code{['L1']} \\
\code{MERA bulk geometry} & Forman mean $-2.242$, Ollivier mean $-0.3297$ \\
artefacts & \code{result/spiral\_v15.png}, \code{result/spiral\_v15\_metrics.png}, \code{result/\_v15\_data.json} \\
wall clock & $\sim 1245$~s; order-of-magnitude reference, machine-dependent \\
\bottomrule
\end{tabular}
\caption{Expected output for a correct reproduction.}
\label{tab:expected}
\end{table}

\begin{quote}
\textbf{Two cautions that belong in this section, not in a footnote.}
\emph{(a)} These numbers certify that the pipeline \emph{ran correctly}. They do not
certify that any conclusion is \emph{true}: $54/54$ means 54 scored assertions hold.
\emph{(b)} The values in Table~\ref{tab:expected} are transcribed from the frozen
reference run, which was re-executed \emph{after} the most recent change to the F6d
guard. The earlier inference --- that renaming a diagnostic and changing its
condition string cannot move the totals, because a diagnostic does not enter the
denominator --- was checked against that re-run and held for all three totals. The
wall-clock figure did move ($1138$~s $\rightarrow$ $1245$~s, $+9.4\%$); it is
machine-dependent and should be read as an order of magnitude only.
\end{quote}

\textbf{Data.} The pipeline reads no external data for the Ising layers; the
Gray--Scott exhibit uses a published dataset, cited in the source.

\textbf{Scope of the accompanying repository.} The code repository also contains
work outside the scope of this paper, including an exploratory framework whose
framing and terminology differ from the account given here. This paper is
self-contained: where the two differ, the claims, numbers, and boundaries stated
here govern, and nothing in the repository's other material should be read as a
result of this paper.

% ---------------------------------------------------------------------------
\section{Conclusion}
\label{sec:conclusion}

We stated an operational criterion for a negative control --- it must be able to
fail when the guarded conclusion is false --- and used it as a filter rather than an
aspiration: it excluded positivity, finiteness, and non-emptiness checks outright,
including one candidate already in place. We sorted the survivors into four types,
built a per-layer ledger, and gave the ledger a completeness guard that checks it
against the \emph{runtime} registry instead of its own text. We separated scored
assertions from diagnostics so that the pass rate cannot be inflated by checks that
cannot fail.

Applied to a seven-layer pipeline, the method produced a ledger covering 6 of 7
layers and, more usefully, three negative results it would have been easy to omit:
a layer whose gap is structural (nothing to push against), a bridge that is fully
wired yet carries no upstream-derived quantity, and a ``geometry'' readout that is
an analytic property of the network wiring rather than an emergent one. It also
caught a real defect in which a completeness criterion was inverted --- a check that
passed for precisely the wrong reason, which is the failure mode this method exists
to find.

The substrate physics is textbook Ising criticality, pinned by a parameter-free CFT
finite-size formula at every accessible $L$ and cross-checked by a genuinely
independent spectral estimator. That is deliberate: the substrate is the
\emph{test bench}. The claim on offer is that the bench is honest --- that the
pipeline reports, in the same document, both what it verified and what it could not.

\medskip
\noindent\textbf{Falsifiability of the method itself.}
The ledger's completeness guard fails if a cited guard does not exist in the
runtime registry or if a layer is missing. The coverage guard fails if coverage
changes. Both are checks on the accounting, which is precisely their scope: they
make the ledger \emph{hard to fake}, not the physics true. This is the method's
most distinctive property: \textbf{the validation framework is itself validated
by the same standard it enforces}. It is falsifiable in exactly the way it
demands of every pipeline guard. The fact that this self-check caught a real
defect --- an inverted completeness criterion that passed for the wrong reason
--- is the strongest evidence that the discipline works, because it worked
\emph{on itself}.

\medskip
\noindent\textbf{Reproducibility statement.} Every numerical claim in this
paper can be regenerated by running \code{python spiral\_model\_v15.py}
from the tagged release at
\url{https://github.com/yanjianzhong/spiral-emergence} (tag \texttt{v15}).
No precomputed data was read; no results were cherry-picked.

% ---------------------------------------------------------------------------
\appendix
\section{Guard accounting, verbatim}
\label{app:accounting}

\noindent Log line, quoted without editing:
\begin{quote}
\code{guards 54/54 = 100.00\% (plus 20 diagnostics excluded, of which 17 did not
pass)}
\end{quote}

\noindent Diagnostic registration is explicit per guard (\code{expect\_pass=False});
74 guards total $=$ 54 scored $+$ 20 diagnostic. Of the 20 diagnostics, 17 fail by
design and 3 pass while remaining excluded from the denominator. The distinction
between ``failed'' and ``defective'' is therefore carried by the data structure and
not by the reader's charity.

\medskip
\noindent\textbf{A methodological note, recorded deliberately.} The suspicion that
the bond dimension constituted an undeclared coupling between layers 4 and 5 ---
described in Sec.~\ref{sec:l4l5} --- was raised and then falsified within seconds,
because the cheapest available falsification test was named \emph{before} any
effort was spent defending the claim. This was the third time in the project that a
causal claim was overturned by its own measurement. The practice of stating the
cheapest falsification first is worth more than any individual result in this paper,
and is recommended independently of it.

\bibliographystyle{unsrt}
\bibliography{spiral_v15}

\end{document}